\subsection{热振子的坐标密度核}
\label{sec:phys-ghs-thermal-mehler}

光子数基底能直接给出热占据；正交分量探测却读取连续的坐标分布。要把这两种描述接起来，需要热密度算符的坐标核及其完整归一化。

先考虑一个正则模态，物理坐标为 $q$，$[q,p]=\ii\hbar$，哈密顿量为
\[
H=\frac{p^2}{2M}+\frac{M\omega^2q^2}{2},\qquad M>0,\ \omega>0.
\]
若 $q$ 是长度，$M$ 是质量；对光场的模态坐标，$M$ 由坐标归一化确定，不能另赋光子的静止质量。相应无量纲正交分量为 $Q=\sqrt{M\omega/\hbar}\,q$、$P=p/\sqrt{M\hbar\omega}$，满足 $[Q,P]=\ii$。

\begin{proposition}[欧氏核及热态]
\label{prop:phys-ghs-mehler-thermal}
定义欧氏时间 $\tau_E>0$ 及 $K_{\tau_E}(q_f,q_i)=\langle q_f|\ee^{-\tau_E H/\hbar}|q_i\rangle$，则
\begin{equation}
\begin{split}
K_{\tau_E}(q_f,q_i)
&=\sqrt{\frac{M\omega}{2\pi\hbar\sinh(\omega\tau_E)}}\\
&\quad\times\exp\left[-\frac{M\omega}{2\hbar\sinh(\omega\tau_E)}
\bigl((q_i^2+q_f^2)\cosh(\omega\tau_E)-2q_iq_f\bigr)\right].
\end{split}
\label{eq:phys-ghs-mehler-kernel}
\end{equation}
对温度 $T>0$，令 $\beta=1/(k_{\mathrm B}T)$、$a=\beta\hbar\omega$，热态坐标核为
\begin{equation}
\rho_\beta(q_f,q_i)=\frac{K_{\beta\hbar}(q_f,q_i)}{Z(\beta)},
\qquad Z(\beta)=\frac1{2\sinh(a/2)}.
\label{eq:phys-ghs-thermal-density-kernel}
\end{equation}
\end{proposition}
\begin{proof}
欧氏作用量为 $I_E=(M/2)\int_0^{\tau_E}(\dot q^2+\omega^2q^2)\dd\tau$。其驻定路径及端点作用量是
\begin{align*}
q_{\mathrm{cl}}(\tau)&=
\frac{\sinh[\omega(\tau_E-\tau)]}{\sinh(\omega\tau_E)}q_i
+\frac{\sinh(\omega\tau)}{\sinh(\omega\tau_E)}q_f,\\
I_E[q_{\mathrm{cl}}]&=\frac M2[q_{\mathrm{cl}}\dot q_{\mathrm{cl}}]_0^{\tau_E}
=\frac{M\omega[(q_i^2+q_f^2)\cosh(\omega\tau_E)-2q_iq_f]}
{2\sinh(\omega\tau_E)}.
\end{align*}
置 $K=C(\tau_E)\ee^{-I_E[q_{\mathrm{cl}}]/\hbar}$，代入热方程
\[
\partial_{\tau_E}K=\frac{\hbar}{2M}\partial_{q_f}^2K
-\frac{M\omega^2q_f^2}{2\hbar}K
\]
并比较常数项，得到 $\partial_{\tau_E}\log C=-(\omega/2)\coth(\omega\tau_E)$。小时间核必须趋于 $\sqrt{M/(2\pi\hbar\tau_E)}\exp[-M(q_f-q_i)^2/(2\hbar\tau_E)]$，唯一确定 \cref{eq:phys-ghs-mehler-kernel} 的前因子。

取对角坐标 $q_f=q_i=q$ 和 $\tau_E=\beta\hbar$，指数化为 $-(M\omega/\hbar)\tanh(a/2)q^2$。普通高斯积分给出
\[
Z=\sqrt{\frac{M\omega}{2\pi\hbar\sinh a}}
\sqrt{\frac{\pi\hbar}{M\omega\tanh(a/2)}}
=\frac1{2\sinh(a/2)}.
\]
所以 $\int\rho_\beta(q,q)\dd q=1$。这一迹包括零点能；删去零点能会改变 $Z$，但在分子分母同步删去时归一化密度算符不变。
\end{proof}

\begin{example}[正交分量涨落与热纯度]
对频率为 $\omega>0$、质量为 $M$ 的谐振子热态，记 $a=\beta\hbar\omega$、$\beta=(k_{\mathrm B}T)^{-1}$。求位置与动量涨落、纯度和熵，并检查低温与高温极限。

\textbf{解。} \label{ex:phys-ghs-thermal-fluctuations}
归一化对角核为
\[
\rho_\beta(q,q)=\sqrt{\frac{M\omega}{\pi\hbar}\tanh(a/2)}
\exp[-(M\omega/\hbar)\tanh(a/2)q^2].
\]
由高斯二阶矩及 $q,p$ 的相空间旋转对称性，
\[
\langle q^2\rangle=\frac{\hbar}{2M\omega}\coth(a/2),\qquad
\langle p^2\rangle=\frac{M\hbar\omega}{2}\coth(a/2).
\]
因此 $\langle Q^2\rangle=\langle P^2\rangle=\bar n+1/2$，$\bar n=(\ee^a-1)^{-1}$；也有 $-\partial_\beta\log Z=(\hbar\omega/2)\coth(a/2)$。低温极限给出基态涨落。高温 $a\ll1$ 时，展开 $\coth(a/2)=2/a+a/6+O(a^3)$，得到
\[
\langle q^2\rangle=\frac{k_{\mathrm B}T}{M\omega^2}
\left[1+\frac{a^2}{12}+O(a^4)\right].
\]

用热半群的拼接性计算纯度，
\[
\Tr\rho_\beta^2=\frac{Z(2\beta)}{Z(\beta)^2}
=\tanh(a/2)=\frac1{2\bar n+1}.
\]
数态概率 $p_n=(1-\ee^{-a})\ee^{-an}$ 则给出
\[
S(\rho_\beta)=-\sum_n p_n\log_2p_n
=(\bar n+1)\log_2(\bar n+1)-\bar n\log_2\bar n.
\]
这与后文从辛本征值算出的高斯态熵相同。这里 $\tau_E=\beta\hbar$ 是热密度核的欧氏时间，不能作为实验中的实时传播时间代入幺正演化。
\end{example}

核的拼接也可以在坐标中直接计算。两段核相乘后，中间坐标 $q$ 的二次系数是 $M\omega[\coth(\omega\tau_1)+\coth(\omega\tau_2)]$；对它完成平方，积分前因子为
\[
\sqrt{\frac{2\pi\hbar}
{M\omega[\coth(\omega\tau_1)+\coth(\omega\tau_2)]}}.
\]
利用 $\coth u+\coth v=\sinh(u+v)/(\sinh u\sinh v)$，指数和前因子共同还原总时间 $\tau_1+\tau_2$ 的核。有限模态可逐模使用这一恒等式；无限模态的行列式乘积还需要收敛或正则化条件。
